\newcommand{\titulus}{\nomenFesti{S. Bonaventuræ, Episcopi et Ecclesiæ Doctoris.}
\dies{Die 15. Iulii.}}
\newcommand{\oratio}{\pars{Oratio.}

\noindent Da, quǽsumus, omnípotens Deus, ut, beáti Bonaventúræ, epíscopi, natalícia celebrántes, et ipsíus proficiámus eruditióne præclára et caritátis ardórem iúgiter æmulémur.

\pars{Pro pace in universo mundo.} \scriptura{Sir. 50, 25; 2 Esdr. 4, 20; \textbf{H416}}

\vspace{-4mm}

\antiphona{II D}{temporalia/ant-dapacemdomine.gtex}

\vfill

\noindent Deus, a quo sancta desidéria, recta consília et iusta sunt ópera: da servis tuis illam, quam mundus dare non potest, pacem; ut et corda nostra mandátis tuis dédita, et hóstium subláta formídine, témpora sint tua protectióne tranquílla.

\noindent Per Dóminum nostrum Iesum Christum, Fílium tuum, qui tecum vivit et regnat in unitáte Spíritus Sancti, Deus, per ómnia sǽcula sæculórum.

\noindent \Rbardot{} Amen.}
\newcommand{\invitatorium}{\pars{Invitatorium.}

\vspace{-4mm}

\antiphona{IV}{temporalia/inv-mirabileminsanctis.gtex}}
\newcommand{\hymnusmatutinum}{\pars{Hymnus}

\cuminitiali{IV}{temporalia/hym-LaetiColentes.gtex}}
\newcommand{\matversus}{\noindent \Vbardot{} Iustum dedúxit Dóminus per vias rectas.

\noindent \Rbardot{} Et osténdit illi regnum Dei.}
\newcommand{\lectioi}{\pars{Lectio I.} \scriptura{1 Reg. 21, 1-14}

\noindent De libro primo Regum.

\noindent Témpore illo vínea erat Naboth Iezrahelítæ, quæ erat in Iézrahel iuxta palátium Achab regis Samaríæ.

\noindent Locútus est ergo Achab ad Naboth dicens: «Da mihi víneam tuam, ut fáciam mihi hortum hólerum, quia vicína est et prope domum meam.

\noindent Dabóque tibi pro ea víneam meliórem aut, si tibi commódius putas, argénti prétium quanto digna est».

\noindent Cui respóndit Naboth: «Propítius mihi sit Dóminus, ne dem hereditátem patrum meórum tibi».

\noindent Venit ergo Achab in domum suam tristis et indígnans super verbo, quod locútus fúerat ad eum Naboth Iezrahelítes dicens: «Non dabo tibi hereditátem patrum meórum». Et proíciens se in léctulum suum avértit fáciem ad paríetem et non comédit panem.

\noindent Ingréssa est autem ad eum Iézabel uxor sua dixítque ei: «Quid est hoc, unde ánima tua contristáta est? Et quare non cómedis panem?».

\noindent Qui respóndit ei: «Quia locútus sum Naboth Iezrahelítæ et dixi ei: Da mihi víneam tuam, accépta pecúnia; aut, si tibi placet, dabo tibi víneam pro ea. Et ille ait: “Non dabo tibi víneam meam”».

\noindent Dixit ergo ad eum Iézabel uxor eius: «Grandis auctoritátis es et bene regis regnum Israel! Surge et cómede panem et æquo esto ánimo: ego dabo tibi víneam Naboth Iezrahelítæ».

\noindent Scripsit ítaque lítteras ex nómine Achab et signávit eas ánulo eius et misit ad maióres natu et ad optimátes, qui erant in civitáte eius et habitábant cum Naboth. Litterárum autem hæc erat senténtia: «Prædicáte ieiúnium et sedére fácite Naboth in cápite pópuli et submíttite duos viros fílios Bélial contra eum, et testimónium dicant: “Maledixísti Deum et regem”; et edúcite eum et lapidáte, sicque moriátur».

\noindent Fecérunt ergo cives eius maióres natu et optimátes, qui habitábant cum eo in urbe, sicut præcéperat eis Iézabel et sicut scriptum erat in lítteris, quas míserat ad eos. {\color{gray} Prædicavérunt ieiúnium et sedére fecérunt Naboth in cápite pópuli; et ingréssi duo viri fílii Bélial sedérunt contra eum et illi, ut viri diabólici, dixérunt contra eum testimónium coram multitúdine: «Maledíxit Naboth Deum et regem». Quam ob rem eduxérunt eum extra civitátem et lapídibus interfecérunt;} miserúntque ad Iézabel dicéntes: «Lapidátus est Naboth et mórtuus est».}
\newcommand{\responsoriumi}{\pars{Responsorium 1.} \scriptura{\Rbardot{} Ps. 37, 7-8.12 \Vbardot{} ibid., 13; \textbf{H165}}

\vspace{-5mm}

\responsorium{IV}{temporalia/resp-totadiecontristatus-CROCHU.gtex}{}}
\newcommand{\lectioii}{\pars{Lectio II.} \scriptura{1 Reg. 21, 15-21.27-29}

\noindent Factum est autem, cum audísset Iézabel lapidátum Naboth et mórtuum, locúta est ad Achab: «Surge, pósside víneam Naboth Iezrahelítæ, qui nóluit tibi acquiéscere et dare eam, accépta pecúnia; non enim vivit Naboth, sed mórtuus est».

\noindent Quod cum audísset Achab, mórtuum vidélicet Naboth, surréxit et descendébat in víneam Naboth Iezrahelítæ, ut possidéret eam.

\noindent Factus est ígitur sermo Dómini ad Elíam Thesbíten dicens: «Surge et descénde in occúrsum Achab regis Israel, qui est in Samaría; ecce est in vínea Naboth, ad quam descéndit, ut possídeat eam.

\noindent Et loquéris ad eum dicens: Hæc dicit Dóminus: Occidísti, ínsuper et possedísti!

\noindent Et post hæc addes: Hæc dicit Dóminus: In loco, in quo linxérunt canes sánguinem Naboth, lambent tuum quoque sánguinem».

\noindent Et ait Achab ad Elíam: «Num invenísti me, inimíce mi?».

\noindent Qui dixit: «Invéni, eo quod venúmdatus sis, ut fáceres malum in conspéctu Dómini.

\noindent Ecce ego indúcam super te malum et démetam posterióra tua et interfíciam de Achab quidquid masculíni sexus sive impúberem sive púberem in Israel».

\noindent Itaque cum audísset Achab sermónes istos, scidit vestem suam et opéruit cilício carnem suam ieiunavítque et dormívit in sacco et ambulábat demísso cápite.

\noindent Factus est autem sermo Dómini ad Elíam Thesbíten dicens: «Nonne vidísti humiliátum Achab coram me?

\noindent Quia ígitur humiliátus est mei causa, non indúcam malum in diébus eius, sed in diébus fílii sui ínferam malum dómui eius».}

\newcommand{\responsoriumii}{\pars{Responsorium 2.} \scriptura{\Rbardot{} Ps. 79, 14.16 \Vbardot{} ibid., 15; \textbf{H88}}

\vspace{-5mm}

\responsorium{III}{temporalia/resp-devastavitvineam-CROCHU.gtex}{}

\vfill

\rubrica{vel ad libitum:}

\vspace{3mm}

\pars{Responsorium 2.} \scriptura{\Rbardot{} Ier. 2, 21 \Vbardot{} ibid.; \textbf{H217}}

\vspace{-5mm}

\responsorium{VIII}{temporalia/resp-vineamea-CROCHU.gtex}{}}
\newcommand{\lectioiii}{\pars{Lectio III.} \scriptura{Cap. 7, 1. 2. 4. 6: Opera omnia, 5, 312-313}

\noindent Ex Opúsculo sancti Bonaventúræ epíscopi De itinerário mentis in Deum.

\noindent Christus est via et óstium. Christus est scala et vehículum tamquam \emph{propitiatórium super arcam Dei collocátum et sacraméntum a sǽculis abscónditum.} Ad quod propitiatórium qui áspicit plena conversióne vultus, aspiciéndo eum in cruce suspénsum, per fidem, spem et caritátem, devotiónem, admiratiónem, exsultatiónem, appretiatiónem, laudem et iubilatiónem; \emph{pascha,} hoc est tránsitum, cum eo facit, ut per virgam crucis tránseat mare Rubrum, ab Ægýpto intrans desértum, ubi gustat manna abscónditum, et cum Christo requiéscat in túmulo quasi extérius mórtuus, séntiens tamen, quantum possíbile est secúndum statum viæ, quod in cruce dictum est latróni cohærénti Christo: \emph{Hódie mecum eris in paradíso.}

\noindent In hoc autem tránsitu, si sit perféctus, opórtet quod relinquántur omnes intellectuáles operatiónes, et apex afféctus totus transferátur et transformétur in Deum. Hoc autem est mýsticum et secretíssimum, quod nemo novit, nisi qui áccipit, nec áccipit nisi qui desíderat, nec desíderat nisi quem ignis Spíritus Sancti medúllitus inflámmat, quem Christus misit in terram. Et ídeo dicit Apóstolus, hanc mýsticam sapiéntiam esse per Spíritum Sanctum revelátam.

\noindent Si autem quæras quómodo hæc fiant, intérroga grátiam, non doctrínam; desidérium, non intelléctum; gémitum oratiónis, non stúdium lectiónis; sponsum, non magístrum; Deum, non hóminem; calíginem, non claritátem; non lucem, sed ignem totáliter inflammántem et in Deum excessívis unctiónibus et ardentíssimis affectiónibus transferéntem. Qui quidem ignis Deus est, et hic camínus est in Ierúsalem, et Christus hunc accéndit in fervóre suæ ardentíssimæ passiónis, quem solus ille vere pércipit, qui dicit: \emph{Suspéndium elégit ánima mea, et mortem ossa mea.} Quam mortem qui díligit, vidére potest Deum, quia indubitánter verum est: \emph{Non vidébit me homo et vivet.} Moriámur ígitur et ingrediámur in calíginem, imponámus siléntium sollicitudínibus, concupiscéntiis et phantasmátibus; transeámus cum Christo crucifíxo \emph{ex hoc mundo ad Patrem,} ut, osténso nobis Patre, dicámus cum Philíppo: \emph{Súfficit nobis}; audiámus cum Paulo: \emph{Súfficit tibi grátia mea}; exsultémus cum David dicéntes: \emph{Déficit caro mea et cor meum, Deus cordis mei et pars mea Deus in ætérnum. Benedíctus Dóminus in ætérnum et dicat omnis pópulus: Fiat, Fiat.}}
\newcommand{\responsoriumiii}{\pars{Responsorium 3.} \scriptura{\Rbardot{} Sir. 15, 3 \Vbardot{} Ps. 131, 18; \textbf{H62}}

\vspace{-5mm}

\responsorium{VII}{temporalia/resp-cibavitillum-CROCHU-cumdox.gtex}{}}
\newcommand{\hymnuslaudes}{\pars{Hymnus}

\cuminitiali{VIII}{temporalia/hym-ORedemptorisPietas.gtex}}
\newcommand{\lectiobrevis}{\pars{Lectio Brevis.} \scriptura{Sap. 7, 13-14}

\noindent Sapiéntiam sine fictióne dídici et sine invídia commúnico; divítias illíus non abscóndo. Infinítus enim thesáurus est homínibus; quem qui acquisiérunt, ad amicítiam in Deum se paravérunt propter disciplínæ dona commendáti.}
\newcommand{\responsoriumbreve}{\pars{Responsorium breve.} \scriptura{Sir. 45, 9}

\antiphona{VI}{temporalia/resp-amaviteum.gtex}}
\newcommand{\preces}{\noindent Christum Deum sanctum, fratres, exaltémus,~\gredagger{} orántes ut serviámus illi in sanctitáte et iustítia coram ipso ómnibus diébus nostris,~\grestar{} et acclamémus:

\Rbardot{} Tu solus sanctus, Dómine.

\noindent Qui tentári voluísti per ómnia pro similitúdine nostra absque peccáto,~\grestar{} miserére nostri, Dómine Iesu.

\Rbardot{} Tu solus sanctus, Dómine.

\noindent Qui nos omnes ad perfectiónem caritátis vocásti,~\grestar{} sanctífica nos, Dómine Iesu.

\Rbardot{} Tu solus sanctus, Dómine.

\noindent Qui nos iussísti esse salem terræ et lucem mundi,~\grestar{} illúmina nos, Dómine Iesu.

\Rbardot{} Tu solus sanctus, Dómine.

\noindent Qui voluísti ministráre,~\gredagger{} non ministrári,~\grestar{} fac nos tibi et frátribus humíliter servíre, Dómine Iesu.

\Rbardot{} Tu solus sanctus, Dómine.

\noindent Tu, splendor glóriæ Patris et figúra substántiæ eius,~\grestar{} fac ut in glória vultum tuum respiciámus, Dómine Iesu.

\Rbardot{} Tu solus sanctus, Dómine.}
\newcommand{\benedictus}{\pars{Canticum Zachariæ.} \scriptura{Eccli. 15, 5; \textbf{H61}}

%\vspace{-4mm}

{
\grechangedim{interwordspacetext}{0.18 cm plus 0.15 cm minus 0.05 cm}{scalable}%
\antiphona{II D}{temporalia/ant-spiritusapientiae.gtex}
\grechangedim{interwordspacetext}{0.22 cm plus 0.15 cm minus 0.05 cm}{scalable}%
}

%\vspace{-3mm}

\scriptura{Lc. 1, 68-79}

%\vspace{-2mm}

\cantusSineNeumas
\initiumpsalmi{temporalia/benedictus-initium-ii-D-auto.gtex}

%\vspace{-1.5mm}

\input{temporalia/benedictus-ii-D.tex} \Abardot{}}
\include{hebdomadaxv}
\include{feriaiv}
